% Options for packages loaded elsewhere
\PassOptionsToPackage{unicode}{hyperref}
\PassOptionsToPackage{hyphens}{url}
\PassOptionsToPackage{dvipsnames,svgnames,x11names}{xcolor}
\documentclass[
  a4paper,
]{article}
\usepackage{xcolor}
\usepackage[margin=25mm]{geometry}
\usepackage{amsmath,amssymb}
\setcounter{secnumdepth}{-\maxdimen} % remove section numbering
\usepackage{iftex}
\ifPDFTeX
  \usepackage[T1]{fontenc}
  \usepackage[utf8]{inputenc}
  \usepackage{textcomp} % provide euro and other symbols
\else % if luatex or xetex
  \usepackage{unicode-math} % this also loads fontspec
  \defaultfontfeatures{Scale=MatchLowercase}
  \defaultfontfeatures[\rmfamily]{Ligatures=TeX,Scale=1}
\fi
\usepackage{lmodern}
\ifPDFTeX\else
  % xetex/luatex font selection
  \setmainfont[]{DejaVu Serif}
  \setsansfont[]{DejaVu Sans}
  \setmonofont[]{DejaVu Sans Mono}
\fi
% Use upquote if available, for straight quotes in verbatim environments
\IfFileExists{upquote.sty}{\usepackage{upquote}}{}
\IfFileExists{microtype.sty}{% use microtype if available
  \usepackage[]{microtype}
  \UseMicrotypeSet[protrusion]{basicmath} % disable protrusion for tt fonts
}{}
\makeatletter
\@ifundefined{KOMAClassName}{% if non-KOMA class
  \IfFileExists{parskip.sty}{%
    \usepackage{parskip}
  }{% else
    \setlength{\parindent}{0pt}
    \setlength{\parskip}{6pt plus 2pt minus 1pt}}
}{% if KOMA class
  \KOMAoptions{parskip=half}}
\makeatother
\usepackage{longtable,booktabs,array}
\newcounter{none} % for unnumbered tables
\usepackage{calc} % for calculating minipage widths
% Correct order of tables after \paragraph or \subparagraph
\usepackage{etoolbox}
\makeatletter
\patchcmd\longtable{\par}{\if@noskipsec\mbox{}\fi\par}{}{}
\makeatother
% Allow footnotes in longtable head/foot
\IfFileExists{footnotehyper.sty}{\usepackage{footnotehyper}}{\usepackage{footnote}}
\makesavenoteenv{longtable}
\usepackage{graphicx}
\makeatletter
\newsavebox\pandoc@box
\newcommand*\pandocbounded[1]{% scales image to fit in text height/width
  \sbox\pandoc@box{#1}%
  \Gscale@div\@tempa{\textheight}{\dimexpr\ht\pandoc@box+\dp\pandoc@box\relax}%
  \Gscale@div\@tempb{\linewidth}{\wd\pandoc@box}%
  \ifdim\@tempb\p@<\@tempa\p@\let\@tempa\@tempb\fi% select the smaller of both
  \ifdim\@tempa\p@<\p@\scalebox{\@tempa}{\usebox\pandoc@box}%
  \else\usebox{\pandoc@box}%
  \fi%
}
% Set default figure placement to htbp
\def\fps@figure{htbp}
\makeatother
\setlength{\emergencystretch}{3em} % prevent overfull lines
\providecommand{\tightlist}{%
  \setlength{\itemsep}{0pt}\setlength{\parskip}{0pt}}
\usepackage{luatexja}
\usepackage{bookmark}
\IfFileExists{xurl.sty}{\usepackage{xurl}}{} % add URL line breaks if available
\urlstyle{same}
\hypersetup{
  pdftitle={The Twelfth Thought Experiment: From the Fine-Structure Constant to Observation Maps and Gauge Symmetry},
  pdfauthor={Noriaki Kihara (WF System Co., Ltd.) ORCID 0009-0004-6753-4020},
  colorlinks=true,
  linkcolor={Maroon},
  filecolor={Maroon},
  citecolor={Blue},
  urlcolor={Blue},
  pdfcreator={LaTeX via pandoc}}

\title{The Twelfth Thought Experiment: From the Fine-Structure Constant
to Observation Maps and Gauge Symmetry}
\usepackage{etoolbox}
\makeatletter
\providecommand{\subtitle}[1]{% add subtitle to \maketitle
  \apptocmd{\@title}{\par {\large #1 \par}}{}{}
}
\makeatother
\subtitle{― Starting from the Fine-Structure Constant, Rereading
Spacetime, Mass, Charge, Color Charge and Gauge Symmetry from the
Light-Cone Equation and the Choice of Observable Axes ―}
\author{Noriaki Kihara (WF System Co., Ltd.) ORCID 0009-0004-6753-4020}
\date{October 8, 2026 Version DOI 10.5281/zenodo.23226259}

\ltjsetparameter{jacharrange={-2,-3}}
\begin{document}
\maketitle

\textbf{Author:} Noriaki Kihara\\
\textbf{ORCID:} 0009-0004-6753-4020\\
\textbf{Version DOI:} 10.5281/zenodo.23226259\\
\textbf{Concept DOI:} 10.5281/zenodo.23226258\\
\textbf{Version:} v1.0\\
\textbf{Date:} 2026-10-08\\
\textbf{Related papers:} Design document (Paper 0) {[}S1{]}, the
Eleventh Thought Experiment {[}S2{]}, nontrivial quadratic closure
{[}S3{]}, reexamination of xyztRQ {[}S4{]}, composition of central
projection {[}S5{]}

\textbf{Keywords:} fine-structure constant, Bohr model, light cone,
Minkowski metric, proper time, quadratic closure, observation map,
indistinguishability, gauge symmetry, Kaluza--Klein, SU(5)

\emph{This is an English translation of the Japanese original
(thought\_experiment\_12\_observation\_mapping\_gauge\_symmetry\_ja\_v1.0.md).}

\begin{center}\rule{0.5\linewidth}{0.5pt}\end{center}

\section{Abstract}\label{abstract}

In the previous paper, the Eleventh Thought Experiment {[}S2{]}, the
fine-structure constant \(\alpha\) was used as a known input value as it
stands. This paper starts by asking where that \(\alpha\) comes from. In
the Bohr ground state of hydrogen \(\alpha=v/c\), so \(\alpha\) can
first be read as a velocity ratio. The Lorentz factor
\(\gamma=1/\sqrt{1-\alpha^2}\) is only a component ratio of
\(r^2+\tau^2+(it)^2=0\), which is the light-cone equation
\(x^2+y^2+z^2+(it)^2=0\) with one proper-time axis added. A worldline
slower than light lies on the light cone in the space with one more
axis.

We extend this light-cone equation to eight directions, the three
spatial directions plus the not directly readable \(R\), \(Q_1\),
\(Q_2\), \(Q_3\): \(x^2+y^2+z^2+R^2+Q_1^2+Q_2^2+Q_3^2+(it)^2=0\). We
give the axes no physical names from the start and treat them as eight
unnamed directions \((u_1,\ldots,u_8)\). When an observer chooses a
directly readable subspace, the remaining directions cannot be
distinguished and appear only as a sum of squares. A transformation
\(G\) for which the observation map \(P\) satisfies \(P(U)=P(GU)\) is an
observational symmetry, so symmetry can be read as the mathematical
expression of indistinguishability. An observer who reads only \(t\)
sees an \(O(7)\)-type symmetry, and if the five hidden directions split
as 3+2,
\(S(U(3)\times U(2))\simeq(SU(3)\times SU(2)\times U(1))/\mathbb Z_6\)
appears as a candidate. This is the same structure as the 3+2 split of
the SU(5) grand unified theory.

This paper is a conceptual note without computation; the eight
dimensions, the signature and the 3+2 split from §4 on are put in by
hand as assumptions.

\begin{center}\rule{0.5\linewidth}{0.5pt}\end{center}

\section{0. Position of this paper}\label{position-of-this-paper}

The Eleventh Thought Experiment {[}S2{]} computed two-body systems of
standard theory numerically, referenced only to the rest frame of an
absorber at infinity. The input \(\alpha\) was the CODATA value {[}1{]},
and its meaning was not asked.

This paper is the work in the opposite direction. It does not accept
\(\alpha\) as an input value; it starts from ``what is \(\alpha\) in the
first place'', returns to the light-cone equation, removes the names of
the axes, and rereads spacetime, mass, charge, color charge and gauge
symmetry from what an observer can distinguish. It is also a rewriting,
in geometric language, of the first-principle candidate of the design
document {[}S1{]}, ``states exist that have no names but are
distinguishable from one another and can be counted'', and of the
position that ``a state itself cannot be read directly''.

This paper is a conceptual note without computation; the eight
dimensions, the signature and the 3+2 split from §4 on are put in by
hand as assumptions.

\begin{center}\rule{0.5\linewidth}{0.5pt}\end{center}

\section{1. Introduction}\label{introduction}

This thought experiment started from a very simple question.

\textbf{Where does the fine-structure constant \(\alpha\) come from in
the first place?}

In the standard treatment it is defined as

\[
\alpha=
\frac{e^2}{4\pi\varepsilon_0\hbar c}
\simeq\frac{1}{137.036}
\]

{[}1{]}, and in quantum electrodynamics it is treated as the
dimensionless coupling constant that expresses the strength of the
electromagnetic interaction.

But if we go back to the historical starting point {[}2, 3, 4{]}, in the
Bohr model of the hydrogen atom the ground state satisfies

\[
\frac{v}{c}=\alpha .
\]

That is,

\[
\boxed{\alpha=\frac{\text{velocity of the electron in the hydrogen ground state}}{\text{speed of light}}}
\]

Therefore

\[
\alpha^{-1}\simeq137
\]

can first be read, naively, as the velocity ratio

\[
c\simeq137\,v .
\]

Starting from this simple fact, might we be able to reread the
fine-structure constant, special relativity, the Minkowski metric, and
even mass, charge, color charge and gauge symmetry from a different
angle?

This paper follows that question step by step.

\begin{center}\rule{0.5\linewidth}{0.5pt}\end{center}

\section{2. Reading the fine-structure constant as a velocity
ratio}\label{reading-the-fine-structure-constant-as-a-velocity-ratio}

In the Bohr model of the hydrogen atom {[}2{]},

\[
\frac{mv^2}{r}
=
\frac{e^2}{4\pi\varepsilon_0r^2}
\]

and

\[
mvr=\hbar
\]

give

\[
v=
\frac{e^2}{4\pi\varepsilon_0\hbar}.
\]

Dividing by the speed of light \(c\),

\[
\frac vc
=
\frac{e^2}{4\pi\varepsilon_0\hbar c}
=\alpha.
\]

Therefore in the hydrogen ground state

\[
\boxed{\frac vc=\alpha}
\]

It is then no surprise that \(\alpha\) appeared when Sommerfeld put
special-relativistic corrections into the hydrogen atom {[}3{]}.

The Lorentz factor of special relativity is

\[
\gamma=
\frac1{\sqrt{1-v^2/c^2}},
\]

so substituting

\[
\frac vc=\alpha
\]

gives

\[
\gamma=
\frac1{\sqrt{1-\alpha^2}}.
\]

In the low-velocity approximation,

\[
\gamma
=
1+\frac12\alpha^2
+\frac38\alpha^4+\cdots .
\]

In this form, the reason why \(\alpha^2,\alpha^4,\ldots\) appear in the
fine structure is very intuitive.

But here an even simpler view is possible.

Why do we have to think from \(\gamma\) at all?

\begin{center}\rule{0.5\linewidth}{0.5pt}\end{center}

\section{3. Seeing it from the light-cone equation rather than the
Lorentz
factor}\label{seeing-it-from-the-light-cone-equation-rather-than-the-lorentz-factor}

Set \(c=1\). The propagation of light is expressed by Minkowski's
light-cone equation

\[
x^2+y^2+z^2-t^2=0 .
\]

Writing time as the imaginary number \(it\) {[}5{]},

\[
\boxed{x^2+y^2+z^2+(it)^2=0}
\]

The Lorentz transformations are determined as the linear transformations
that preserve this form. What comes first is the light-cone equation;
the Lorentz transformations appear afterwards as the transformations
that preserve it.

Next consider a particle moving with velocity \(v\). With the distance
traveled \(r\) and the proper time \(\tau\),

\[
t^2-r^2=\tau^2 .
\]

Moving \(\tau^2\) to the left-hand side gives

\[
\boxed{r^2+\tau^2+(it)^2=0}
\]

This has the same form as the light-cone equation. A worldline that is
slower than light in \((r,t)\) lies on the light cone in \((r,\tau,t)\),
where one proper-time axis \(\tau\) has been added {[}6, 11{]} (Fig.
1(a)(b)).

Dividing by \(t^2\),

\[
\left(\frac rt\right)^2+
\left(\frac \tau t\right)^2=1.
\]

Reading here

\[
\frac rt=\frac vc=\alpha ,
\]

we have

\[
\boxed{
\left(\frac \tau t\right)^2+\alpha^2=1
}
\]

and

\[
\frac \tau t=\sqrt{1-\alpha^2}.
\]

Taking the reciprocal,

\[
\gamma=\frac t\tau
=\frac1{\sqrt{1-\alpha^2}} .
\]

Therefore \(\gamma\) is not fundamental.

What comes first is

\[
\boxed{r^2+\tau^2+(it)^2=0},
\]

the light-cone equation with one axis added, and the Lorentz factor is
nothing but its component ratio read from one axis.

In this paper we call a zero sum of squares of this form, containing one
imaginary component, a \textbf{quadratic closure}. The name belongs to
this paper, but the equation is the light-cone equation itself. That a
zero sum of squares can have a nontrivial solution only if at least one
component is not real was shown in {[}S3{]}.

Here the way the problem looks changes.

\begin{center}\rule{0.5\linewidth}{0.5pt}\end{center}

\section{4. Extending spacetime}\label{extending-spacetime}

Next, using ordinary three-dimensional space and splitting the
\(\tau^2\) of §3 into two directions \(R\) and \(Q\), we put

\[
\boxed{
x^2+y^2+z^2+R^2+Q^2=t^2
}
\]

Moving terms,

\[
x^2+y^2+z^2
+R^2+Q^2+(it)^2=0.
\]

This is also the light-cone equation, only with more axes. The imaginary
axis is \(t\) alone.

What matters here is not to fix the physical meaning of each axis from
the start.

First, simply, we read

\[
(x,y,z)
\]

as the directions that can be measured directly as space, and

\[
(R,Q)
\]

as the directions that cannot be measured directly in the same way.

Then

\[
t^2-(x^2+y^2+z^2)
=
R^2+Q^2
\]

can be read as the observational relation

\[
\boxed{
\begin{array}{c}
\text{difference of squares of time and the directly readable part}\\
=\ \text{sum of squares of the not directly readable part}
\end{array}
}
\]

For an observer who reads \(xyz\), this equation has the same form as
the equation of central projection. When a point \(P=(x,y,z)\) on the
tangent plane at distance \(R_c\) from the center \(O\) is mapped to the
sphere of radius \(R_c\) by \(\pi(P)=R_c\,P/\ell\) {[}S5{]}, the
distance \(\ell\) from \(O\) to \(P\) satisfies

\[
\ell^2=x^2+y^2+z^2+R_c^2 .
\]

Comparing with the equation above, the curvature radius in the
projection direction is the square root of the sum of squares of the not
directly readable directions, \(R_c=\sqrt{R^2+Q^2}\), and the distance
\(\ell\) from the center to the point corresponds to \(t\). The proper
time \(\tau\) of §3 corresponds to \(R_c\), the spatial distance \(r\)
to the distance \(|x|\) on the tangent plane, and the slope \(v=r/t\) of
the worldline coincides with \(\sin\theta=|x|/\ell\) of the angle
\(\theta\) of the projection ray. Removing the imaginary unit \(i\) from
the light-cone equation gives the Pythagorean equation of central
projection (Fig. 1(c)).

\pandocbounded{\includegraphics[keepaspectratio,alt={Figure 1}]{figs/fig01_lightcone_central_projection.png}}
\textbf{Figure 1.} The light cone and central projection are the same
right triangle. (a) The light cone in the \((r,t)\) plane. A worldline
of slope \(v<1\) lies inside the cone. (b) Adding the proper-time axis
\(\tau\) gives \(r^2+\tau^2+(it)^2=0\), and the same worldline lies on
the surface of the cone. In the cut at time \(t\), a right triangle
forms with legs \(r\) (blue) and \(\tau\) (green) and the cut radius
\(t\) (red) as hypotenuse. (c) Central projection
\(\ell^2=|x|^2+R_c^2\). The same triangle appears as the distance
\(|x|\) on the tangent plane (blue), the curvature radius \(R_c\) in the
projection direction (green), and the distance \(\ell\) from the center
to the point (red). \(r\leftrightarrow|x|\),
\(\tau=\sqrt{R^2+\sum Q_i^2}\leftrightarrow R_c\),
\(t\leftrightarrow\ell\), \(\sin\theta=v\). A schematic drawn for a
generic \(v\); there are no data.

At this point we do not need to assume that

\begin{itemize}
\tightlist
\item
  \(t\) is essentially a ``time axis''
\item
  \(R\) is essentially a ``mass axis''
\item
  \(Q\) is essentially a ``charge axis''.
\end{itemize}

These may be names attached afterwards to observational results. The
names \(t,R,Q\) are those used in {[}S4{]}. The reading that relates the
fifth coordinate to rest mass also appears in Wesson's Space-Time-Matter
theory {[}12{]}.

\begin{center}\rule{0.5\linewidth}{0.5pt}\end{center}

\section{5. Perhaps there is only one kind of
momentum}\label{perhaps-there-is-only-one-kind-of-momentum}

Suppose a state is expressed by one phase

\[
\Phi .
\]

Then the momentum in each direction can in general be written as its
phase gradient,

\[
K_A=\partial_A\Phi .
\]

For the ordinary spatial directions we read it as

\[
(p_x,p_y,p_z).
\]

Applying the same idea to the other axes, there is only one gradient
vector

\[
K=
(K_x,K_y,K_z,K_t,K_R,K_Q).
\]

It is possible that what we call

\[
K_t\rightarrow E,
\]

\[
K_R\rightarrow m,
\]

\[
K_Q\rightarrow q
\]

are its components.

Applying the closure of §4 to this gradient gives

\[
K_t^2=K_x^2+K_y^2+K_z^2+K_R^2+K_Q^2,
\]

that is,

\[
E^2=p^2+m^2+q^2,
\]

which returns to the ordinary mass shell \(E^2=p^2+m^2\) for \(q=0\).
The reading in which the momentum in a hidden direction is read as
charge and its square is added to the mass shell is the same structure
as in Kaluza--Klein theory {[}7, 8{]}, where a particle moving on the
light cone in five dimensions carries its momentum in the hidden
direction as charge in four dimensions and its square as the
four-dimensional mass.

But here too the order has to be reversed.

What really exists from the start is not

\[
E,m,q .
\]

There are first unnamed gradient components, and we may merely be giving
the names

\[
\text{energy},
\quad
\text{mass},
\quad
\text{charge}
\]

to the ways they appear in observation.

Then the view

\[
\boxed{
\begin{array}{c}
\text{momentum, energy, mass and charge are}\\
\text{projections of one gradient onto different directions}
\end{array}
}
\]

becomes possible.

\begin{center}\rule{0.5\linewidth}{0.5pt}\end{center}

\section{6. Splitting Q into three
directions}\label{splitting-q-into-three-directions}

There is no need to take \(Q\) as a single axis either.

Rather, we put the three directions

\[
(Q_1,Q_2,Q_3).
\]

The whole is then the eight-dimensional

\[
\boxed{
(x,y,z,t,R,Q_1,Q_2,Q_3)
}
\]

The quadratic closure is

\[
\boxed{
x^2+y^2+z^2
+R^2+Q_1^2+Q_2^2+Q_3^2
=
t^2
}
\]

or

\[
\boxed{
x^2+y^2+z^2
+R^2+Q_1^2+Q_2^2+Q_3^2+(it)^2=0
}
\]

It is the light-cone equation in eight dimensions. For an observer who
reads \(xyz\), the curvature radius of the central projection is
\(R_c=\sqrt{R^2+Q_1^2+Q_2^2+Q_3^2}\).

It has been shown in the Kaluza--Klein framework that obtaining the
Standard Model group as isometries of the hidden directions requires at
least seven directions outside the four spacetime directions {[}13{]}.
In this paper the directions outside the four spacetime directions are
the four \(R,Q_1,Q_2,Q_3\), which is not enough as isometries. This
paper reads the group not from isometries but from indistinguishability
(§9--§11). A precedent that reads the Standard Model group from eight
dimensions (the octonions) is Furey {[}14{]}.

Here too we do not think that the physical meanings

\[
x,y,z,t,R,Q_1,Q_2,Q_3
\]

exist from the start.

Originally we think there are only

\[
(u_1,u_2,\ldots,u_8),
\]

\textbf{eight unnamed directions}.

When an observer can read three of them directly as space, it names them

\[
(u_1,u_2,u_3)\rightarrow(x,y,z).
\]

It may merely be naming the features observed in the remaining
directions

\[
t,\quad R,\quad Q_1,Q_2,Q_3 .
\]

Here the meaning of physical quantities itself becomes an
observer-dependent map.

\begin{center}\rule{0.5\linewidth}{0.5pt}\end{center}

\section{7. Observers who can read different
axes}\label{observers-who-can-read-different-axes}

From here a strong thought experiment becomes possible.

That we directly read

\[
(x,y,z)
\]

does not mean that this is the only possible observation map.

For example, an observer that can directly read only

\[
(Q_1,Q_2,Q_3)
\]

can be defined mathematically.

For that observer,

\[
(Q_1,Q_2,Q_3)
\]

looks like ordinary three-dimensional space.

Conversely,

\[
x,y,z,t,R
\]

become internal directions that cannot be read directly.

The physical phenomena of that world would look extremely different from
those of ours.

But if the state and the interactions on the eight-dimensional side are
the same, this is nothing but

\[
\boxed{
\text{different observation maps of the same physics}
}
\]

This is not ``different physical laws'' but

\[
\boxed{
\text{the result of choosing a different subspace as the observed space}
}
\]

\begin{center}\rule{0.5\linewidth}{0.5pt}\end{center}

\section{8. What is seen if only the t axis can be
read}\label{what-is-seen-if-only-the-t-axis-can-be-read}

Consider an even more extreme case.

Consider an observer that can directly read, of the eight directions,
only

\[
t .
\]

The remaining seven directions,

\[
x,y,z,R,Q_1,Q_2,Q_3 ,
\]

cannot be distinguished individually.

What can then be observed is not their breakdown but only the sum of
squares

\[
\boxed{
\rho^2=
x^2+y^2+z^2
+R^2
+Q_1^2+Q_2^2+Q_3^2
}
\]

That is, it appears only as

\[
\boxed{
t^2=\rho^2
}
\]

To this observer, the eight-dimensional light cone is seen only as its
radius \(\rho=t\).

A complicated state with seven degrees of freedom degenerates into one
scalar quantity.

This is important.

The unobservable degrees of freedom have not disappeared.

\textbf{Because they cannot be distinguished, they appear only as one
norm.}

In general, if only \(k\) of the eight dimensions can be read directly,

\[
\mathbb R^8
\rightarrow
\mathbb R^k_{\rm observed}
\oplus
\mathbb R^{8-k}_{\rm hidden} .
\]

If the hidden side cannot be distinguished individually, it is observed
as one quantity

\[
\rho_{\rm hidden}^2
=
\sum_{\rm hidden}u_a^2 .
\]

Therefore

\[
\boxed{
\begin{array}{c}
\text{dimensional reduction is not that dimensions}\\
\text{physically disappear, but that the observer}\\
\text{can no longer distinguish them}
\end{array}
}
\]

can be read.

\begin{center}\rule{0.5\linewidth}{0.5pt}\end{center}

\section{9. Symmetry appears from
here}\label{symmetry-appears-from-here}

This thought experiment gives another view of why standard theory
requires strong symmetries.

Let the observation map be

\[
P .
\]

If, after acting with a transformation \(G\) on a state \(U\),

\[
P(U)=P(GU),
\]

then the observer cannot distinguish the two.

Therefore

\[
\boxed{
P(U)=P(GU)
\quad\Longrightarrow\quad
G\text{ is an observational symmetry transformation}
}
\]

That is,

\[
\boxed{
\text{symmetry is the mathematical expression of indistinguishability}
}
\]

can be read. The connection between gauge symmetry and observation and
relations is discussed by Rovelli {[}15{]}.

There is no need to impose symmetry on nature afterwards.

If exchanging or rotating axes that the observer cannot name leaves the
observational results unchanged, that is directly a symmetry.

\begin{center}\rule{0.5\linewidth}{0.5pt}\end{center}

\section{10. The huge symmetry when only t can be
read}\label{the-huge-symmetry-when-only-t-can-be-read}

When only \(t\) can be read, only

\[
\rho^2=
x^2+y^2+z^2+R^2+Q_1^2+Q_2^2+Q_3^2
\]

can be observed.

Considering only the real squared norm, rotating the seven hidden
directions leaves \(\rho\) unchanged.

Therefore a large symmetry of type

\[
O(7)
\]

appears.

What matters is that

\[
\boxed{
\begin{array}{c}
\text{the fewer the observable dimensions,}\\
\text{the larger the indistinguishability on the hidden side,}\\
\text{and the larger the apparent symmetry}
\end{array}
}
\]

(Fig. 2).

\pandocbounded{\includegraphics[keepaspectratio,alt={Figure 2}]{figs/fig02_observers_symmetry.png}}
\textbf{Figure 2.} Three observation maps of the same eight directions
and the size of the rotational symmetry of the hidden directions. Left:
the directions that can be read directly (blue) and the directions that
cannot be distinguished (gray) for an observer reading \(xyz\), an
observer reading \(Q_1Q_2Q_3\) and an observer reading only \(t\).
Right: the number \(k\) of directly readable directions and the
dimension \((8-k)(7-k)/2\) of the rotation group of the hidden
directions. Reading only \(t\) gives \(O(7)\) (21 dimensions); reading
\(xyz\) gives 10 dimensions. When \(t\) is hidden the group becomes
\(O(7-k,1)\), but the dimension is the same.

Conversely, when the observational resolution rises and directions that
were indistinguishable become distinguishable, the symmetry becomes
smaller.

Therefore the interpretation

\[
\boxed{
\begin{array}{c}
\text{symmetry breaking}
=
\text{directions that could not be distinguished}\\
\text{becoming distinguishable}
\end{array}
}
\]

also becomes possible.

\begin{center}\rule{0.5\linewidth}{0.5pt}\end{center}

\section{11. The 3+2 split}\label{the-32-split}

Suppose we are observers who directly read

\[
(x,y,z).
\]

Then on the hidden side the five directions

\[
(t,R,Q_1,Q_2,Q_3)
\]

remain.

Suppose that observation makes these five directions distinguishable
only as the two groups

\[
(Q_1,Q_2,Q_3)
\]

and

\[
(t,R),
\]

that is, as

\[
\boxed{3+2}
\]

If the three internal components, as complex states, are
indistinguishable under change of basis, a symmetry of type

\[
U(3)
\]

becomes a candidate.

Likewise, for the two components a symmetry of type

\[
U(2)
\]

becomes a candidate.

Therefore the structure

\[
U(3)\times U(2)
\]

can appear naturally.

Further, imposing a common overall phase / determinant condition and
considering

\[
S(U(3)\times U(2)),
\]

we have

\[
S(U(3)\times U(2))
\simeq
\frac{
SU(3)\times SU(2)\times U(1)
}{
\mathbb Z_6
}
\]

{[}10{]}.

This coincides with the global form
\((SU(3)\times SU(2)\times U(1))/\mathbb Z_6\) usually adopted as the
gauge group of the Standard Model.

This 3+2 split is the same structure as the one in which
\(S(U(3)\times U(2))\subset SU(5)\) remains in the SU(5) grand unified
theory {[}9, 10{]} (Fig. 3).

\pandocbounded{\includegraphics[keepaspectratio,alt={Figure 3}]{figs/fig03_split_3_2.png}}
\textbf{Figure 3.} The 3+2 split of the five hidden directions. Mixing
within \((Q_1,Q_2,Q_3)\) (blue) corresponds to \(U(3)\), mixing within
\((t,R)\) (orange) to \(U(2)\), and mixing between the two groups (gray)
can be distinguished by observation. Imposing a common overall phase /
determinant condition gives \(S(U(3)\times U(2))\subset SU(5)\),
\(S(U(3)\times U(2))\simeq(SU(3)\times SU(2)\times U(1))/\mathbb Z_6\)
{[}9, 10{]}.

However, at this stage this is \textbf{not a derivation}.

In particular, the quadratic closure taken as basic in this research,

\[
\sum z_n^2=0 ,
\]

and the Hermitian norm preserved by the ordinary unitary groups,

\[
\sum |z_n|^2 ,
\]

are not the same.

Obtaining \(U(n)\) from eight real directions requires newly assuming a
complex structure, and this paper does not derive it.

Therefore, to claim

\[
3+2
\rightarrow
SU(3)\times SU(2)\times U(1),
\]

it will be necessary to make clear which observable is taken as the
invariant.

Still, it is noteworthy that from the observational structure

\[
\boxed{
\text{3 indistinguishable directions}
+
\text{2 indistinguishable directions}
}
\]

alone, a mathematical structure extremely close to the Standard Model
gauge group appears as a candidate.

\begin{center}\rule{0.5\linewidth}{0.5pt}\end{center}

\section{12. Reversing the viewpoint}\label{reversing-the-viewpoint}

In the usual description,

\[
\text{space},
\quad
\text{time},
\quad
\text{mass},
\quad
\text{charge},
\quad
\text{color charge}
\]

are introduced from the start as different kinds of physical quantities.

Then symmetry is required to connect them.

In this thought experiment the order is reversed.

What exists first is only the unnamed state space

\[
\boxed{
(u_1,\ldots,u_8)
}
\]

The observer chooses from it a subspace it can distinguish directly.

The rest cannot be distinguished directly.

And the results of projecting them onto the observable space are named

\[
\text{time},
\quad
\text{mass},
\quad
\text{charge},
\quad
\text{color charge}.
\]

Then

\[
\boxed{
\begin{array}{c}
\text{a physical quantity is not an attribute proper to an axis}\\
\text{but a reading name that depends on the observation map}
\end{array}
}
\]

Likewise,

\[
\boxed{
\begin{array}{c}
\text{symmetry is not a constraint added to nature}\\
\text{but the existence of degrees of freedom the observer cannot distinguish}
\end{array}
}
\]

\begin{center}\rule{0.5\linewidth}{0.5pt}\end{center}

\section{13. Back to the fine-structure
constant}\label{back-to-the-fine-structure-constant}

Here we return to the fine-structure constant, our starting point.

In the hydrogen ground state,

\[
\alpha=\frac vc.
\]

At first this looked like a mere velocity ratio.

But in the view of this thought experiment,

\[
v
\]

and

\[
c
\]

may also be ratios of reading an unnamed higher-dimensional gradient
from different observation directions.

Therefore the possibility arises that it can be interpreted as

\[
\boxed{
\alpha
=
\text{ratio of gradient components onto two observation directions}
}
\]

Then the fine-structure constant may not be merely ``the strength of the
electromagnetic interaction'' but

\[
\boxed{
\begin{array}{c}
\text{an angle or ratio arising when a higher-dimensional state}\\
\text{is projected onto a particular observed subspace}
\end{array}
}
\]

At this stage it is undetermined which of

\[
Q_i/t,
\quad
Q_i/R,
\quad
|\mathbf Q|/R,
\quad
\text{or another projection angle}
\]

it corresponds to.

But at least, starting from the historical fact

\[
\alpha=\frac vc ,
\]

we place a starting point for looking for \(\alpha\) as a purely
geometric ratio.

\begin{center}\rule{0.5\linewidth}{0.5pt}\end{center}

\section{14. Verification items from this thought
experiment}\label{verification-items-from-this-thought-experiment}

At least the following verifiable directions arise.

First,

\[
\boxed{
\begin{array}{c}
\text{from the same eight-dimensional dynamics,}\\
\text{different effective physics should be obtained}\\
\text{only by changing the observed subspace}
\end{array}
}
\]

Therefore, without changing the interaction law at all, the effective
motion can be computed numerically for each case, such as reading

\[
xyz ,
\]

reading

\[
Q_1Q_2Q_3 ,
\]

or reading only

\[
t .
\]

Second,

\[
\boxed{
\begin{array}{c}
\text{the fewer the observable dimensions,}\\
\text{the larger the effective symmetry should be}
\end{array}
}
\]

Third,

\[
\boxed{
\begin{array}{c}
\text{when the observational resolution rises and}\\
\text{previously indistinguishable internal directions are separated,}\\
\text{the symmetry should break}
\end{array}
}
\]

Fourth, at least some of mass, charge and color charge may be
reproducible as

\[
\boxed{
\begin{array}{c}
\text{not different basic attributes but}\\
\text{different projected components of the same unnamed state space}
\end{array}
}
\]

Fifth, it can be tested whether the Standard Model symmetry of type

\[
SU(3)\times SU(2)\times U(1)
\]

can arise from

\[
\boxed{
\begin{array}{c}
\text{the observational indistinguishability classes}\\
\text{of 3 directions and 2 directions}
\end{array}
}
\]

\begin{center}\rule{0.5\linewidth}{0.5pt}\end{center}

\section{15. Conclusion}\label{conclusion}

This thought experiment started by questioning the meaning of one
dimensionless constant,

\[
\boxed{\alpha\simeq1/137}
\]

In the hydrogen ground state,

\[
\alpha=\frac vc.
\]

Reading this not from the Lorentz factor but from

\[
r^2+\tau^2+(it)^2=0 ,
\]

the light-cone equation with one axis added, \(\alpha\) appears as a
simple direction ratio.

Extending this to

\[
x^2+y^2+z^2+R^2+Q_1^2+Q_2^2+Q_3^2
=
t^2
\]

gives the eight-dimensional structure

\[
(x,y,z,t,R,Q_1,Q_2,Q_3).
\]

But there is no need to give these eight axes physical names from the
start.

Essentially there exists only the unnamed state space

\[
(u_1,\ldots,u_8),
\]

and it can be thought that physical meanings such as

\[
\text{space},
\quad
\text{time},
\quad
\text{mass},
\quad
\text{charge},
\quad
\text{color}
\]

arise according to

\[
\boxed{
\text{which axes are read as the observable space}
}
\]

Then

\[
\boxed{
\text{different physics}
=
\text{different basic laws}
}
\]

need not hold.

Rather,

\[
\boxed{
\text{different physics}
=
\text{different observation maps of the same unnamed physics}
}
\]

may hold.

Further, if we read

\[
\boxed{
\text{symmetry}
=
\text{transformations the observer cannot distinguish}
}
\]

the reason why standard theory requires strong gauge symmetries can also
be understood in another form.

And symmetry breaking may be

\[
\boxed{
\begin{array}{c}
\text{directions that could not be distinguished}\\
\text{becoming distinguishable}
\end{array}
}
\]

In this view, time, mass, charge, color charge and gauge symmetry are
not independent concepts from the start.

They may all be grasped in a unified way as

\[
\boxed{
\begin{array}{c}
\text{the results of reading one unnamed higher-dimensional state}\\
\text{from a limited observed subspace}
\end{array}
}
\]

\begin{center}\rule{0.5\linewidth}{0.5pt}\end{center}

\section{The last question}\label{the-last-question}

If the physical law on the eight-dimensional side were only one, closed
by

\[
\boxed{
\sum_{a=1}^{8} z_a^2=0
}
\]

alone, then is the very fact that we call

\[
xyz
\]

``space'',

\[
t
\]

``time'',

\[
R
\]

``mass and curvature'', and

\[
Q_1,Q_2,Q_3
\]

``charge and color charge''

\textbf{a physical law, or merely the observer's choice?}

\begin{center}\rule{0.5\linewidth}{0.5pt}\end{center}

\section{Reproducing the figures}\label{reproducing-the-figures}

Figures 1--3 are schematics without data. Running
\texttt{figures/make\_figures.py} (Python 3.9, numpy, matplotlib)
generates SVG and PNG files with English labels in the same folder.

{\def\LTcaptype{none} % do not increment counter
\begin{longtable}[]{@{}ll@{}}
\toprule\noalign{}
Fig. & File \\
\midrule\noalign{}
\endhead
\bottomrule\noalign{}
\endlastfoot
1 &
\texttt{figures/fig01\_lightcone\_central\_projection.\{svg,png\}} \\
2 & \texttt{figures/fig02\_observers\_symmetry.\{svg,png\}} \\
3 & \texttt{figures/fig03\_split\_3\_2.\{svg,png\}} \\
\end{longtable}
}

\begin{center}\rule{0.5\linewidth}{0.5pt}\end{center}

\section{Record of AI involvement}\label{record-of-ai-involvement}

\begin{itemize}
\tightlist
\item
  \textbf{ChatGPT} (2026-10-08): first draft of this paper (structure
  and text of §1--§15 and ``The last question'').
\item
  \textbf{Claude} (2026-10-08): correction of escapes produced by export
  from Google Docs, rewriting §3 to start from the light-cone equation,
  unifying the signs of §4, §6 and §15 to a single imaginary axis \(t\),
  the correspondence between the mass shell and Kaluza--Klein in §5, the
  statements on the global form, SU(5) and complex structure in §11,
  addition of the abstract, §0 and references with bibliographic checks
  (\texttt{引用文献台帳\_20261008.md}), the correspondence with central
  projection in §4, Figures 1--3 and the script
  \texttt{figures/make\_figures.py}, and this English translation.
  Design and judgment are the author's.
\end{itemize}

\begin{center}\rule{0.5\linewidth}{0.5pt}\end{center}

\section{References}\label{references}

\subsection{This research series}\label{this-research-series}

{[}S1{]} N. Kihara, \textbf{Design document (Paper 0)}. Concept DOI:
10.5281/zenodo.22851944.

{[}S2{]} N. Kihara, \textbf{The Eleventh Thought Experiment: Emission
Records of Two-Body Systems without a Background Spacetime - Computing
Standard Theory (Dirac-Coulomb Bound States, 1PN Two-Body Gravity,
Einstein A Coefficients) Referenced Only to the Rest Frame of an
Absorber at Infinity, and Reading Out Emission-Absorption Relations for
the Hydrogen Atom, Electron-Electron, Proton-Proton and the Hydrogen
Molecule}, v1.0 (2026-10-07). Concept DOI: 10.5281/zenodo.23199299;
Version DOI: 10.5281/zenodo.23199300.

{[}S3{]} N. Kihara, \textbf{Complex Phase Structure Emerging from
Nontrivial Quadratic Closure: A Foundational Connection to Finite-Order
Exchange Systems and Discrete Born-Type Weights}, v1.0 (2026-07-18).
Concept DOI: 10.5281/zenodo.21422505; Version DOI:
10.5281/zenodo.21422506.

{[}S4{]} N. Kihara, \textbf{Reexamination of the 6-Dimensional Sign
Vector xyztRQ: A Thought Experiment Toward a Future Paper}, v4
(2026-04-30). Concept DOI: 10.5281/zenodo.19902677; Version DOI:
10.5281/zenodo.19904714.

{[}S5{]} N. Kihara, \textbf{Composition of Central Projection and the
Closed Form of the Composite Curvature Radius: An Algebraic Formulation
of High-Dimensional Reduction via One Central Projection and Commutative
Cuts on the Sphere}, v1 (2026-05-07). Concept DOI:
10.5281/zenodo.20060728; Version DOI: 10.5281/zenodo.20060729.

\subsection{External references}\label{external-references}

{[}1{]} P. J. Mohr, D. B. Newell, B. N. Taylor, E. Tiesinga, CODATA
recommended values of the fundamental physical constants: 2022,
Rev.~Mod. Phys. 97, 025002 (2025).

{[}2{]} N. Bohr, On the constitution of atoms and molecules, Phil. Mag.
26, 1 (1913).

{[}3{]} A. Sommerfeld, Zur Quantentheorie der Spektrallinien, Ann. Phys.
51, 1 (1916).

{[}4{]} H. Kragh, Magic number: A partial history of the fine-structure
constant, Arch. Hist. Exact Sci. 57, 395 (2003).

{[}5{]} H. Minkowski, Raum und Zeit, Phys. Z. 10, 104 (1909).

{[}6{]} L. C. Epstein, \emph{Relativity Visualized} (Insight Press, San
Francisco, 1983).

{[}7{]} Th. Kaluza, Zum Unitätsproblem der Physik, Sitzungsber. Preuss.
Akad. Wiss. Berlin (Math. Phys.), 966 (1921).

{[}8{]} O. Klein, Quantentheorie und fünfdimensionale
Relativitätstheorie, Z. Phys. 37, 895 (1926).

{[}9{]} H. Georgi, S. L. Glashow, Unity of all elementary-particle
forces, Phys. Rev.~Lett. 32, 438 (1974).

{[}10{]} J. Baez, J. Huerta, The algebra of grand unified theories,
Bull. Amer. Math. Soc. 47, 483 (2010).

{[}11{]} J. M. C. Montanus, Proper-time formulation of relativistic
dynamics, Found. Phys. 31, 1357 (2001).

{[}12{]} P. S. Wesson, \emph{Space-Time-Matter: Modern Kaluza-Klein
Theory} (World Scientific, Singapore, 1999).

{[}13{]} E. Witten, Search for a realistic Kaluza-Klein theory, Nucl.
Phys. B 186, 412 (1981).

{[}14{]} C. Furey, SU(3)\_C × SU(2)\_L × U(1)\_Y (× U(1)\_X) as a
symmetry of division algebraic ladder operators, Eur. Phys. J. C 78, 375
(2018).

{[}15{]} C. Rovelli, Why gauge?, Found. Phys. 44, 91 (2014).

\end{document}
